% Part: lambda-calculus
% Chapter: introduction
% Section: basic-pr-lambda

\documentclass[../../../include/open-logic-section]{subfiles}

\begin{document}

\olfileid{lam}{int}{bas}
\olsection{基本原始递归函数是$\lambda$-可定义的}

\begin{lem}
函数 $\Zero$、$\Succ$ 和 $\Proj{n}{i}$ 都是 $\lambd$-可定义的。
\end{lem}

\begin{proof}
$\Zero$ 就是
$\lambd[x][\lambd[y][y]]$。

后继函数~$\Succ$ 由
$\fn{Succ}(u) = \lambd[x][\lambd[y][x(uxy)]]$ 定义。请思考为何该定义可行：对于每个被视为迭代器的数码 $\num{n}$ 和每个函数~$f$，$\fn{Succ}(\num{n},f)$ 是这样一个函数，它在以~$y$ 为输入时，先从~$y$ 起将 $f$ 应用 $n$ 次，然后再多应用一次。

至于投影函数则是不言自明的：$\fn{Proj}^n_i(x_0, \dots,
x_{n-1}) = x_i$。换言之，根据我们的约定，$\fn{Proj}^n_i$ 就是 λ-项 $\lambd[x_0][\dots \lambd[x_{n-1}][x_i]]$。
\end{proof}

\end{document}